\chapter{总结与展望}\label{chap:conclusion}

\section{全文总结}

本文面向 AOSP 15 的 LoongArch 平台 ART，围绕执行引擎适配、正确性验证和性能优化展开研究。研究表明：首先，AOSP 15 的 ART 执行引擎能够在 LoongArch 平台上顺利完成适配，构建出可运行且可验证的执行环境；其次，将目标机回归、libcore 测试与统一数据集整合至同一实验设计，是确保性能对比具备充分可信度的前提；此外，在探索的多项优化方案中，LoongArch64 LSX 浮点 SIMD 向量化带来了稳定且显著的性能提升，构成了本研究中证据链最完整的核心结论。

为此，本文主要完成了以下工作：补齐了汇编器、C++ 解释器、Nterp、JIT 编译器后端、JNI 编译器及若干运行时快速路径，为 LoongArch 平台建立了完整的执行能力；搭建了覆盖编译、部署、目标机回归、libcore 验证和基准测试的实验框架，并将单次对比、重复测量、启动期案例与补充材料进行了合理的验证分层；基于固定对照版本，对三项核心优化方案进行了统一评测，并根据性能提升表现区分了方案的适用场景与收益强度。

实验结果显示，方案一对 \texttt{DaCapo/pmd} 等短生命周期工作负载具有明显的加速效果，但存在一定的测量波动；方案二针对字符串热路径进行了改进，适用于字符串密集场景下的定向微调；方案三在 \texttt{SPECjvm2008/scimark.lu.small} 上取得了稳定且显著的性能增长，是本文最为核心的性能优化成果。

\begin{figure}[htbp]
    \centering
    \includegraphics[width=0.96\textwidth]{conclusion/conclusion-layered-summary.pdf}
    \caption{本文从执行能力构建到主结论凝练的层级总结}
    \label{fig:conclusion-layered-summary}
\end{figure}

如图~\ref{fig:conclusion-layered-summary} 所示，本文的研究逻辑呈现层层递进的结构：由建立基础执行能力起步，经过严谨的实验设计与多维度的案例验证（涵盖控制流、数据流与启动期路径），最终收束至基于 LoongArch64 最小可用 LSX 浮点 SIMD 路径的核心优化结论。该结构体现了从基础系统适配到定向局部优化、再到核心端到端结论的系统性推演过程。

\section{主要结论与贡献凝练}

\subsection{建立可运行的 LoongArch 平台 ART 执行环境}

本文以 AOSP 15 ART 为基础，围绕执行引擎核心模块，完成了 LoongArch 平台上的解释器、Nterp、JIT 编译器后端、JNI 编译器和运行时入口点等适配工作，构建了可运行、可验证的新平台 ART 执行环境。与较早版本 ART 或单纯 JVM 移植相比，AOSP 15 的执行引擎涉及更多子系统，模块耦合也更为紧密。为此，本文以执行路径为主线推进适配工作，避免了零散模块开发的割裂感，从而为后续在国产处理器平台上持续完善 Android 运行时提供了扎实的工程基础。

\subsection{建立支撑性能结论的实验设计}

本文建立了覆盖目标机 run-test、libcore 测试和基准测试评测的实验设计，并将正式单次对比与同日重复测量明确分离。此举旨在将“系统正确性”与“性能变化”严格解耦，防止因功能异常、环境漂移或设备状态波动造成性能收益的误判。该实验框架下的验证层次、标准数据集与有效性威胁控制手段，共同构成了本文性能优化的证据基础。

\subsection{提出基于工作负载特征的优化筛选机制}

面对繁杂的性能热点，本文未采用无差别的优化策略，而是围绕工作负载主路径、收益强度及机制可解释性进行了严格的评估筛选。在此基础上，提取出针对控制流、数据流与启动期路径的代表性优化案例，进而推导出最具应用潜力的核心结论。这一机制确保了优化方案紧密契合实际研究问题，有效避免了边缘冗余方向对系统性能主线的干扰。

\subsection{得到面向代表性工作负载的优化结论}

在固定实验设置下，本文最终保留三项正式优化方案。方案一通过调整基线 JIT 的进入时机，在 \texttt{pmd} 上展现出一定的性能提升趋势，但受测量波动影响，需对其普适性持谨慎态度；方案二针对字符串热路径进行改进，在 \texttt{pmd} 上获得小幅性能增长，属于面向字符串密集型负载的局部优化；方案三补齐了 LoongArch64 最小可用的 LSX 浮点 SIMD 路径，在 \texttt{lu.small} 上取得了 \texttt{+44.40\%} 的单次对比收益，并在重复测量中稳定保持在 \texttt{+45.70\%} 的增幅，构成了本研究最为显著的优化成果。

\subsection{明确 LoongArch 平台 ART 优化的适用边界}

除单方案收益外，本文还明确了执行引擎优化的两项核心边界。首先，性能优化表现出显著的工作负载相关性，不同优化策略所覆盖的执行路径各异，局部收益难以简单外推为通用结论；其次，单独有效的局部优化策略在组合使用时往往无法直接叠加收益，实验表明上述三项方案在叠加状态下并未产生协同增效。这表明，LoongArch 平台 ART 的优化必须紧扣具体工作负载的主导成本展开，只有当优化机制精确覆盖高频路径且额外开销可控时，底层改动才能切实转化为端到端的性能收益。

\section{后续工作展望}

\subsection{完善 AOT 与安装期编译支持}

后续可继续围绕 \texttt{dex2oat} 和安装期编译流程，完善 LoongArch 平台在 AOT 路径上的支持，进一步提升冷启动与安装期执行效率。与本文当前聚焦的 JIT 执行阶段不同，AOT 路径更直接影响安装期成本和冷启动早期表现。若能将安装期 profile、JIT 热点迁移和 LoongArch 后端代码质量优化结合起来，ART 在该平台上的整体性能上限仍有继续提升的空间\citep{AndroidRuntime2024,ArtJitCompiler2024}。

\subsection{扩展 LoongArch 向量化能力}

后续工作可在本文已验证的最小 LSX 浮点 SIMD 指令降级（lowering）基础上，进一步结合 LoongArch 的 LSX 与 LASX 向量扩展能力，面向更多算子、数据类型及复杂循环模式开展优化。补充的消融实验表明，虽然更完整的 SIMD 覆盖在部分特定工作负载上存在收益，但在全部 \texttt{JIT} 负载的几何平均性能上并未超越当前正式方案三所采用的最小可用版本。因此，未来的重点不仅在于扩充降级支持，更需深入探究哪些向量化模式适合默认开启，并探索如何将这些模式与高层循环识别、惯用语（idiom）匹配及启发式触发条件相结合，从而确保新增能力仅在具备实际优势的工作负载中生效\citep{LoongArchISAv2,LoongArchToolchain2025}。

\subsection{扩展基准测试与应用级验证}

未来可在现有实验框架的基础上，引入更为丰富的稳定基准测试集及真实的复杂应用场景，如 SPECjvm 的其他数值工作负载、DaCapo 的扩展子项，以及具备现代 JVM 并发特征的测试集。通过将测试范围逐步拓展至更多样化的工作负载，可更全面地评估 LoongArch 平台 ART 在复杂场景下的性能边界，并进一步明确特定优化策略是应设为默认启用、按需配置，还是仅作为实验性选项\citep{Prokopec2019Renaissance}。

\subsection{增强运行时观测与因果分析}

后续工作还可继续围绕热点预算分析、工作负载分类建模和运行时统计接口建设，建立更系统的 LoongArch 平台 ART 执行引擎优化方法。特别是针对 JIT 分层时机、调试信息的固定成本以及代码缓存利用率等问题，若能依托运行时环境提供更为稳定、细粒度的统计支持，将大幅提升后续性能优化研究的解释力与可迁移性\citep{Luo2022Hubble}。

\section{本章小结}

本章对全文工作进行了系统总结。本文完成了 AOSP 15 ART 在 LoongArch 平台上的执行引擎适配，搭建了支撑性能验证的实验框架，并在统一的对照条件下完成了三项核心优化方案的评测。其中，针对 LoongArch64 LSX 浮点 SIMD 的向量化优化在 \texttt{lu.small} 上取得了显著且稳定的性能收益；而 JIT 分层热度阈值调整与运行时字符串搬运批量化优化，则分别探索了短生命周期工作负载与字符串密集场景下的性能改进空间。总体而言，面向 LoongArch 平台 ART 的底层优化具有切实可行性，但相关性能结论必须建立在严谨的实验控制与明确的工作负载边界之上。